% !PRINTED 118
\weberpage{118}
%<<CONT
an element $b$, then the complex $Pb$ is completely identical with $P$. The two sequences
\[
A = a_0,\ a_1,\ a_2,\ \ldots,\ a_{n-1}
\]
\[
Ab = a_0b,\ a_1b,\ a_2b,\ \ldots,\ a_{n-1}b
\]
can therefore differ from one another only in their arrangement. The
passage from $A$ to $Ab$ is thus a permutation of $n$ figures
$0, 1, 2 \ldots n - 1$, and to every element $b$ of $P$ there corresponds such
a permutation, which we shall denote by $\pi_b$. To two different
elements $b, c$ there always correspond two different
permutations $\pi_b, \pi_c$, and to the unit element there corresponds the
identical permutation. If we set
\[
\pi_b = (A,\ Ab), \quad \pi_c = (A,\ Ac),
\]
then we can also denote $\pi_c$ by $(Ab, Abc)$, and from this
it follows that
\weq{2}{\pi_b \pi_c = \pi_{bc},}
i.\,e. the permutations $\pi$ form a permutation group isomorphic with $P$. This permutation group is transitive, since, for example,
the element $a_0$ can pass into any other element of $P$ whatever. We thus have the theorem:

1.\ \gesperrt{Every group of degree $n$ is isomorphic with a
transitive permutation group of $n$ figures.}

There are of course still other permutation groups isomorphic
with a given group, and it would be of particular
interest to form such a group with the smallest possible number of figures.
This problem cannot at present be solved in general.
We must content ourselves here with a few theorems.

We assume that $R$ is any divisor of $P$ of
index $j$, and denote the elements of $R$ by $c$; further, decomposing
$P$ into a system of subsidiary groups, we set
\weq{3}{P = R + Rb_1 + Rb_2 + \cdots + Rb_{j-1}.}
If then $a$ is any element of $P$, the two
systems
\weq{4}{\begin{aligned}
B &= R,\ Rb_1,\ Rb_2,\ \ldots,\ Rb_{j-1}\\
Ba &= Ra,\ Rb_1a,\ Rb_2a,\ \ldots,\ Rb_{j-1}a,
\end{aligned}}
agree, apart from their order, and so to every element $a$ there corresponds a definite permutation $\pi_a$ of the
$j$ figures $0, 1 \ldots j - 1$, which we denote by
\weq{5}{\pi_a = (B,\ Ba)}
%CONT>>
